\section{Composing crushing proofs across visible diagram factors}
\label{sec:factored-crushing}

A visible connected sum can be recognized by constructing and searching
the exteriors of its diagram factors separately.  The new optional entry
point \path{recognize_factored_by_crushing} makes this composition available
directly to the certified crushing backend.  It combines the established
interlacement factorizer with the full diagram-bound recognition certificates
of Section~\ref{sec:normal-crushing-search}, shares work allowances across
children, and supplies an independent verifier for the resulting verdict.

The high-level \path{recognize} router already uses visible factorization,
as described in \path{research_updates.tex}.  The contribution here is the
explicit proof composition and resource accounting at the direct crushing
interface.  Neither that router nor the existing
\path{recognize_by_crushing} entry point changes its defaults.  The wrapper
does not construct a normal cone for the whole connected sum, transport
normal rays between factors, or infer an inequality for the resolving
parameter of such a combined source.

\subsection{The diagram decomposition and independently reconstructed children}

Let $D$ be a validated spherical one-component PD with $n$ crossings.
Traverse its knot component and record the visited crossing labels.  The
result is a cyclic double-occurrence word $W$ of length $2n$.  Two labels
are interlaced when their two occurrences alternate.  The existing
factorization certificate gives an ordered partition
$C_1,\ldots,C_k$ of the crossings and a spanning forest of interlaced
pairs, with one tree on each block.  Labels within blocks and the blocks'
least labels are in increasing order.

The new verifier validates $D$ afresh and obtains its actual traversal.
It checks literal integer labels, the complete partition, each alternating
forest edge, and connectivity within every claimed block.  A separate
nesting-stack scan checks that distinct blocks do not interlace.  Thus the
blocks are exactly the interlacement components; a forest proves their
connectivity without requiring the possibly quadratic interlacement graph
to be constructed or the discovery algorithm to be rerun.

For clarity, the topological step can be seen directly in the projection.
The occurrences of a component with no interlacement with another lie in
a gap between successive occurrences of that other component.  Connectivity
forces this choice of gap to be consistent throughout the component.
Consequently component intervals are nested or disjoint.  An innermost
block occupies a closed interval of the traversal: every crossing in it
is visited twice, and only its two endpoint edges meet the rest of the
projection.  In the projection sphere those two edges give a visible
splitting curve.  Lifting this curve above and below the projection gives
a sphere meeting the knot in two points.  Remove the innermost block and
continue.  This is the visible-splitting argument used by the existing
factorizer; it establishes
\[
 K(D)=K(D_1)\mathbin{\#}\cdots\mathbin{\#}K(D_k).
\]
It concerns these particular visible factors, which need not be prime.

Authenticating the crossing blocks does not authenticate arbitrary child
PD arrays.  The verifier reconstructs each child from the original
traversal, independently of the factor producer.  Filter that traversal
to the darts whose crossing belongs to $C_i$.  Put $n_i=|C_i|$, retain
the crossings in their original increasing order, and let the $j$th
filtered visit enter slot $s$ of one such crossing.  Assign edge label
$j$ to this slot and $(j+1)\bmod 2n_i$ to its opposite slot
$s+2\bmod4$.  All original cyclic crossing slots, including the
under/over distinction, are preserved.  The reconstructed child is
validated as a spherical one-component PD and must equal the packet's
child PD literally.  No test based only on crossing number, writhe, or
a knot polynomial substitutes for this equality.

There are two interface conventions.  With a single interlacement block,
the original normalized PD is retained literally, including its edge
labels.  The crossing-free circle has no crossing blocks but is represented
by one empty child, rather than by an empty list of recognition obligations.
For every nonempty input,
\[
 n_i\ge1,\qquad \sum_{i=1}^k n_i=n,\qquad k\le n.
\]
The verifier never calls the factor discovery routine, the exterior
constructor, or a normal-surface search backend.

\subsection{One nontrivial child or all trivial children}

The exact outer schema is \path{diagram-factor-crushing-v1}, with fields
\path{schema}, \path{status}, \path{input_pd}, \path{factorization},
\path{factor_pds}, and \path{proofs}.  A proof record consists of an
original \path{factor_index} and a complete child crushing-recognition
certificate.  Reordering discovery does not rename the factor indices.

For \path{UNKNOT}, the packet must contain exactly one child proof for
each index $0,\ldots,k-1$, in that order, and every child proof must have
status \path{UNKNOT}.  For \path{KNOTTED}, it must contain exactly one
in-range child index and a verified child certificate with status
\path{KNOTTED}.  In both cases the child verifier receives the independently
reconstructed PD at that index.  A certificate for the full input or for
a different labelled child cannot fill this obligation.  Duplicate indices,
missing children, and boolean values used as integer indices are rejected.

\begin{proposition}[Visible-factor proof composition]
\label{prop:factored-crushing-soundness}
The independent factorwise verifier accepts only correct knot verdicts
for the original PD.  Its entry point is
\path{verify_factored_crushing_recognition}.
An incomplete child proof cannot substitute for a required complete verdict.
\end{proposition}
\begin{proof}
The checked partition and independently reconstructed children establish
the displayed connected-sum decomposition.  Each retained child verdict
then follows from its own source-authenticated crushing-recognition proof.

For two knots separated by a decomposing sphere, their connected sum is
trivial if and only if both factors are trivial.  One direction follows
by joining spanning discs for the two trivial factors.  Conversely, let
the sum bound an embedded disc and make that disc transverse to the
decomposing sphere.  Their intersection consists of an arc joining the
two knot intersections and some circles.  Every circle cuts off a sphere
disc disjoint from the intersection arc; choose one innermost among these
discs, so its interior misses the spanning disc.  Replace the subdisc it
bounds in the spanning disc by this sphere disc, and push slightly off
the sphere.  The replaced subdisc also misses the intersection arc, whose
endpoints lie on the boundary of the spanning disc.  This surgery preserves
an embedded spanning disc and removes at least one intersection circle.
After the circles are removed, cutting the spanning disc along its single
intersection arc gives a spanning disc for each closed factor tangle.
Both factors are therefore trivial.  Induction proves the assertion for
all visible factors.

It follows that all verified trivial children prove \path{UNKNOT}, whereas
one verified nontrivial child proves \path{KNOTTED}.  An inconclusive child,
a partial \path{REDUCED} source chain, or a supplied-source Euler answer
does not meet either required child-certificate condition.
\end{proof}

The producer retains every attempted child result separately in
\path{factor_results}, including inconclusive results and any earlier
successful children.  These records explain work already spent.  The
conclusive outer proof contains only the proofs needed by the proposition:
one nontrivial child or all trivial children.  The two roles must be kept
distinct when counting cost: a small final certificate does not erase the
discovery work of earlier children.

\subsection{Scheduling and globally shared allowances}

The wrapper obtains option defaults from the actual signature of
\path{recognize_by_crushing}, and validates every supported option before
factor work starts.  Children are visited in increasing order of crossing
count, with the original factor index breaking ties.  Each is called once
with the remaining global allowances.  An inconclusive child does not
immediately stop the wrapper; a later child may still conclude within the
remaining allowances.  A verified nontrivial child stops recognition at
once.  Otherwise every factor must conclude trivial before the wrapper
can return \path{UNKNOT}.

This finite schedule makes no fairness or competitiveness claim.  An early
expensive child can consume a shared allowance needed by a later child.
With all operational caps disabled, completeness of each child recognizer
and finiteness of the factor list imply completeness of the wrapper.
With finite caps, a missing required child verdict produces
\path{INCONCLUSIVE}, not a weaker recognition certificate.  Caller callback
exceptions propagate with their identity; they are not converted into
ordinary local resource exhaustion.

The nine shared accounting units are as follows.  A statistic unused by a
selected child method contributes zero.
\begin{center}\small
\begin{tabular}{ll}
Allowance & Charged child statistics\\
\hline
\path{max_sectors} & \path{sectors_attempted}\\
\path{max_bases} & \path{bases_attempted}\\
\path{max_crushes} & \path{crushing_steps} plus \path{boundary_unlayerings}\\
\path{max_pivots} & \path{lp_pivots}\\
\path{max_anchors} & \path{anchors_tested}\\
\path{max_tree_nodes} & \path{nodes_attempted}\\
\path{max_hitting_subsets} & \path{hitting_subsets_tested}\\
\path{max_ellipsoid_iterations} & \path{lp_ellipsoid_iterations}\\
\path{max_feasibility_calls} & \path{lp_feasibility_calls}
\end{tabular}
\end{center}

For an allowance $B_a$, let $u_{ja}$ be the corresponding charged count
returned by the $j$th attempted child.  Before child $r$ begins, pass
\[
 B_a^{(r)}=B_a-\sum_{j<r}u_{ja};
\]
an unbounded allowance remains unbounded.  The child enforces
$0\le u_{ra}\le B_a^{(r)}$, and the wrapper checks the updated total.
Induction gives $\sum_{j\le r}u_{ja}\le B_a$ for every returned prefix.
Counts from inconclusive children and restarted inner searches are included.
For the conflict candidate policy, candidate counts retain their existing
meaning: visited candidate-selection states, including repeated depth
passes, rather than only successful or distinct final resolving sets.

The \path{max_crushes} row includes both strict reduction operations.
Charging normal crushes alone would fail to enforce the shared allowance
when optional inverse boundary layering is enabled.  Ordinary additive
statistics are summed across children; declared peak statistics such as
maximum tableau size and maximum search depth use their maximum.
Initial-source and shelled-source tetrahedron counts are likewise summed
over the attempted factors, rather than reported as the size of a combined
triangulation.  No child result is memoized or reused for another factor.
Polynomial input preparation and proof replay retain the raw driver's existing scope
outside its charged search allowances.  A zero search budget is therefore
not a claim of zero total computation.

\subsection{An additive cost identity and the largest child-stage parameter}

Let $I$ be the indices actually attempted, and let $P\subseteq I$ be the
indices retained in the final proof.  For an inconclusive wrapper call take
$P=\varnothing$.  Write $W_i$ for all work in the attempted child call,
including its inner preparation, search, retries and proof checks; let
$V_i$ be the outer replay cost of its retained proof.  Up to fixed interface
overhead, the work decomposes as
\[
 W(D)=W_{\rm factor}(D)+\sum_{i\in I}W_i
      +W_{\rm outer\ factor\ replay}(D)+\sum_{i\in P}V_i
      +W_{\rm accounting/output}.
\]
The outer factor-replay term is zero when no conclusive packet is produced.
This is an additive decomposition of actual work.  In particular, stopping
on one nontrivial child removes future calls, but does not replace the
preceding sum by the cost of that child alone.  Retaining or serializing
all attempted results likewise charges their actual combined output size.

Factor discovery, child reconstruction and partition replay use
$O(n\log(n+1))$ word operations with $O(\log(n+1))$-bit indices, and hence
polynomial bit work.  For child $i$, its native initial exterior has
$t_{i0}=O(n_i+1)$ tetrahedra by Section~\ref{sec:diagram-exterior}.
Every normal crush or certified inverse layer strictly decreases this
count.  Thus it performs at most $t_{i0}$ strict reductions and at most
$t_{i0}+1$ Euler search stages.  Summing over the children gives $O(n+1)$
stages, with the extra constant covering the crossing-free convention.

\begin{proposition}[Parameterized additive crushing bound]
\label{prop:factored-crushing-cost}
Disable operational caps, enable vertex reduction, and choose the
polynomial LP policy with conflict-directed transversal search.  Let
$T_{is}$ be the authenticated one-vertex source at Euler stage $s$ of
attempted child $i$.  Write $t_{is}=|T_{is}|$, $h_{is}=h(T_{is})$, and
let $\delta_{is}$ be the conflict-size parameter of
Section~\ref{sec:conflict-candidates}.  There is a fixed polynomial $p$
such that discovery, proof construction and replay have bit cost at most
\[
 p(n+1)+\sum_{i\in I}\sum_s
       p(n_i+1)(3\delta_{is})^{h_{is}}.
\]
In particular, with $H=\max(\{0\}\cup\{h_{is}\})$, this is bounded by
\[
 \operatorname{poly}(n+1)\,(c(n+1))^H
\]
for an absolute constant $c$.  The subset candidate policy has the
corresponding bound
$\operatorname{poly}(n+1)(c(n+1))^{O(H)}$.
\end{proposition}
\begin{proof}
Vertex reduction and any optional inverse layers contribute polynomial
work and strict tetrahedron descent.  At each remaining Euler stage,
vertex reduction guarantees a finite one-vertex source.  It has at most
$4t_{is}$ anchor choices, and its source-derived exact linear data have
bit length polynomial in $t_{is}$.  Proposition~\ref{prop:conflict-candidate-bound}
therefore gives
$p(n_i+1)(3\delta_{is})^{h_{is}}$ work at this stage, after enlarging a
single fixed polynomial to include source algebra, positive-ray extraction,
crushing and independently checked proof data.  The polynomial number of
preprocessing stages and the polynomial factor interface fit the first term
or the polynomial factors of the sum.  Outer replay visits only retained
child proofs and has the same additive form.

Now $\delta_{is}\le\max(1,t_{is})=O(n_i+1)\le O(n+1)$ and the total
number of stages is $O(n+1)$.  Monotonicity of a sufficiently large
polynomial absorbs their sum, yielding the displayed bound in $H$.
The subset-policy conclusion follows instead from
Theorem~\ref{thm:transversal-learning-bound} at each stage.
\end{proof}

Consequently a uniform bound
$h_{is}=O(\log^d(n+2))$ for some fixed $d$ throughout all required child
stages would imply quasi-polynomial recognition on that class of visible
factorizations.  This is a conditional statement: it does not supply that
structural hypothesis.  The maximum concerns every authenticated successor,
not just the first source of a child or one observed conflict set.

\subsection{A restricted-family bound from visible factor size}

The complete ternary-tree policy gives a second consequence without any
assumption on $h$.  It is a corollary of the existing decomposition and
tree policies, rather than a new bound for unrestricted knot diagrams.

\begin{corollary}[Small visible factors]
\label{cor:factored-crushing-small-factors}
Use \path{sector_method='tree'}, \path{tree_model='source'},
\path{lp_policy='polynomial'} and \path{reduce_vertices=True}, with all
nine operational caps set to \path{None}.  Put
$s=\max_i n_i$, with $s=0$ for the crossing-free input.  The wrapper
recognizes the input exactly with bit work and complete proof size bounded by
\[
 \operatorname{poly}(n+1)\,2^{O(s)}.
\]
Consequently it is polynomial for fixed $s$, and also for
$s=O(\log(n+2))$.  For every fixed $d$, it is quasi-polynomial on the
class with $s=O(\log^d(n+2))$.
\end{corollary}
\begin{proof}
The native exterior construction has exactly $20\max(1,n_i)$ tetrahedra;
initial shellings and all later reduction operations can only decrease
this number.  Thus $t_{i0}\le20\max(1,n_i)$, and child $i$ has at most
$t_{i0}+1$ Euler stages.  At a stage with $t$ tetrahedra, each branch makes
one ambiguous tetrahedron singleton.  The tree therefore has depth at
most $t$ and at most
\[
 1+3+\cdots+3^t=\frac{3^{t+1}-1}{2}
\]
nodes.  The one-vertex source has at most $4t$ anchor objectives per node,
and at most one additional envelope call.  The polynomial exact LP policy
and the source-algebra bounds of Section~\ref{sec:euler-branch-tree} give
polynomial bit work and polynomial proof data per node.  Positive extraction,
crushing, inverse layering if enabled, and independent replay contribute
polynomial work per reduction or per recorded node.

For one fixed polynomial $p$, enlarging it to include the polynomial number
of child stages gives the additive bound
\[
 p(n+1)+\sum_i p(n_i+1)\,3^{20\max(1,n_i)}.
\]
Here the sum may include every child even if a negative child stops the
actual run early.  Since $\sum_i n_i=n$ and $k\le n$ for $n>0$, this is
at most $\operatorname{poly}(n+1)2^{O(s)}$.  The crossing-free convention
contributes a single constant-size child.  The same bound covers the
encoded proof bytes and all attempted result records: every normal vector,
local rational witness and source-transition packet has polynomial length,
and the number of recorded nodes and stages satisfies the same bounds.
Substituting $s=O(\log^d(n+2))$ gives the stated restricted-family claim.
\end{proof}

This factor-size benefit was already available through the high-level
factor router.  The new interface supplies composed direct-crushing proofs
and shared budgets.  The stated corollary uses the explicit tree policy;
it is not inferred from whichever policy the raw driver's defaults select.

Neither cost statement compares $h(T)$ for an exterior built directly from
the entire input with the parameters of the separately constructed child
exteriors.  No such comparison is needed: the wrapper searches the latter
sources.  An interlacement-connected diagram can still represent a composite
knot with no visible split, and then this interface has one full-size child.
The operation provides neither hidden prime decomposition nor a general
quasi-polynomial unknot-recognition theorem.

\subsection{Implementation checks}

The retained command and output in
\path{data/factored-crushing-test-command.txt} and
\path{data/factored-crushing-tests.txt} record $238$ passing tests with
Regina available and no dependency skips.  The new producer and verifier
contribute $21$ test methods.  One reconstruction test covers $313$
validated spherical diagrams generated from double-occurrence words with
at most four crossings and their crossing-slot choices.  Other checks cover
the exact positive and negative proof requirements, altered partitions and
child PDs, shared allowances including failed work, producer-disabled
independent replay, and callback exception identity.

As a small complete accounting control, the two-curl unknot has two visible
children.  Both are recognized by one native crush each without any LP calls.
With a global allowance of one crush, the same input returns one complete
trivial child and one independently replayable partial child, and the wrapper
remains inconclusive.  These correctness controls establish neither a timing
advantage nor a general bound on visible factor size or resolving parameters.
